% sup1_00_setting: Appendix on Ba3 Prop. 1, part 1, section opening and setting.
% Written proofs of the first five sketched steps of Ba3 Prop. 1 (slices, external-field kernels, the two-domain
% propagator, the summation by parts, the induction (2.16)) at block factor L = 5. Input from a08_balaban_prop1.tex.

\subsection{The sketched steps of Proposition~1, I: slices, external-field kernels, the summation by parts and the induction}\label{sup1:sec}

\Cref{app:a06_BALtorus} reads the perturbative estimate \cite[Prop.~1, (1.33)]{Ba3} at block factor $L=5$ and uses it in the form
\[
|E(G_{\mathrm{ren}},\Lambda,\Phi)|\ \le\ O(1)\,e_k^{d_v}\lambda_k^{d_s}\,e^{-\delta_0d(\Lambda)}\prod\norm{h\Phi}_{1,\alpha_0},
\]
for a localised renormalised class $G_{\mathrm{ren}}$ with localisation $\Lambda=\{\square(v)\}$, $d(\Lambda)$ the length of a shortest tree through the cubes of $\Lambda$. Its proof on \cite[pp.~420--445]{Ba3} consists of twenty sketched steps and seven further sketched places. This appendix writes them out. We number the sketched steps S1--S21 in the order in which they are invoked in the printed proof (S12 is the printed cube split of \cite[p.~433]{Ba3}, which needs no further argument), and the seven further places F1--F7. This part contains the first five:

\begin{center}
\begin{tabular}{lll}
step & \cite{Ba3} page & printed object\\\hline
S1 & 414 & the kernel \cite[(1.16)]{Ba3}, $[G_k(\Omega,\tilde B)V_k]^nG_k(\Omega,\tilde A+\tilde B)[V_kG_k(\Omega,\tilde B)]^{n'}$, as an external field\\
S2 & 414, 424 & $\delta G_k(\Omega,\Omega_2,\tilde B)=G_k(\Omega,\tilde B)-G_k(\Omega_2,\tilde B)$ as an external field: \cite[(2.5)]{Ba3}, first alternative\\
S3 & 424, 426 & the slice decomposition \cite[(2.6)]{Ba3} and the slice bounds \cite[(2.10)--(2.12)]{Ba3}\\
S4 & 425 & the summation by parts \cite[(2.8)--(2.9)]{Ba3} at a vertex of the graph \cite[(2.4)]{Ba3}\\
S5 & 428 & the induction \cite[(2.16)]{Ba3}
\end{tabular}
\end{center}

S1 and S2 rest on S3, so the order below is S3 (\cref{sup1:ssec:S3}), two common tools (\cref{sup1:ssec:WK}), S1 (\cref{sup1:ssec:S1}), S2 (\cref{sup1:ssec:S2}), S4 (\cref{sup1:ssec:S4}) and S5 (\cref{sup1:ssec:S5}). Everything is at $L=5$, $d=3$, on the scale-$k$ lattice with $k\ge1$, in Ba{\l}aban's own model (no shell mass).

\subsubsection{Setting and inputs}\label{sup1:ssec:setting}

\paragraph{Units and scales.} $\eta\defeq L^{-k}$ with $L=5$ and $k\ge1$, so $\eta\le\frac15$; $t\defeq t_k=L^k\eps$; $e_k\defeq e\,t^{1/2}$ is the running charge \cite[(2.23)]{Ba1}. The slice scales are $s_j\defeq L^j\eta=L^{j-k}$, $0\le j\le k-1$, so $\eta=s_0\le s_j\le\frac1L$. Scalar fields take values in $\R^N$; $U(A_b)=\exp(qe_k\eta A_b)$ with $q$ antisymmetric, $\norm q\le1$ \cite[p.~605]{Ba1}; all $U$ commute and are orthogonal. The covariant derivative is $(D_{A,\mu}\phi)(x)=\eta^{-1}(U(A_{x,\mu})\phi(x+\eta e_\mu)-\phi(x))$ \cite[(1.7)]{Ba1}; inner products carry the weight $\eta^d$ and kernels are taken with respect to $\eta^d$. $|\cdot|_\infty$ is the torus sup metric \cite[(1.3)]{Ba1}, used for every distance, depth and exponential; $|\cdot|$ is the Euclidean torus norm, used only in the H\"older quotients of \cite[(1.32)]{Ba3}; $|\cdot|_\infty\le|\cdot|\le\sqrt3|\cdot|_\infty$.

\paragraph{Blocks and domains.} Unit blocks $B^k(y)$ have side $1$; large blocks have side $M$ \cite[p.~607]{Ba1}. $\Omega=B^k(\Omega^{(k)})$ with $\Omega^{(k)}$ a union of large blocks; $\Omega_1$, $\Omega_2$ are of the same kind \cite[p.~414]{Ba3}. The depth of $x$ in $X$ is $\rho_X(x)\defeq\dist_\infty(x,X^c)$.

\paragraph{Averages and propagators.} $(Q_k(A)\phi)(y)=\sum_{x\in B^k(y)}\eta^dU(A(\Gamma^{(k)}_{y,x}))\phi(x)$ along the composite contours of \cite[(2.1)--(2.2)]{Ba1}; $P_k=Q_k^*Q_k$; $H_X(A)\defeq-\Delta^{\eta,N}_{A,X}+m^2t^2+a_kP_k(A)$ and $G_k(X,A)\defeq H_X(A)^{-1}$ \cite[(2.22)]{Ba1}, where $X$ is a union of unit blocks with Neumann conditions or the torus $T_\eta$. We assume $0\le m^2t^2\le c_m$. The vector field is the case $N=d$, $A=0$, $X=T_\eta$, with $\lmu$ in place of $m$ \cite[p.~608]{Ba1}. $a_j=a(1-L^{-2})/(1-L^{-2j})\in[\frac{24}{25}a,a]$ \cite[(2.15)]{Ba1}; put $a_*\defeq\frac{24}{25}a$.

\paragraph{The hypotheses of \cite[p.~412]{Ba3}.} $|\tilde A|$, $|d^\eta\tilde A|$, $|d^\eta\tilde B|$ and their H\"older norms of exponent $\alpha_0$ are $\le c_1p(L^k\eps)$, and $\dist(\supp\tilde A,\partial\Omega)>2r(L^k\eps)$. Put
\[
\sigma\defeq e_k(\norm{\tilde A}_\infty+\norm{d\tilde A}_\infty)\le2c_1e_kp(t),\qquad \theta(\tilde B)\defeq e_k\norm{d\tilde B}_\infty\le c_1e_kp(t),\qquad\theta(\tilde A+\tilde B)\le2c_1e_kp(t).
\]
\emph{Smallness.} We assume throughout
\begin{equation}\label{sup1:eq:E}
2c_1e_k^{1-\beta}p(t)\le\min\{c,1\},\tag{E}
\end{equation}
with $\beta\in(0,1)$ and $c$ the constants of \cite[(2.23)]{Ba1}. It holds for $t\le t_0$, since $e_k^{1-\beta}p(t)=e^{1-\beta}t^{(1-\beta)/2}b_0(1+\ln t^{-1})^{\bar p}\to0$; it is the ``$e(L^k\eps)$ sufficiently small'' of \cite[Prop.~2.1]{Ba1}.

\paragraph{Depth hypotheses.} The lemmas below are stated with explicit depth hypotheses on the supports of the vertex localisations and of $\tilde A$ (S1: $\supp\tilde A$ at depth $\ge R_0+2$ and vertex supports at depth $\ge R_0+3$ in $\Omega$; S2: vertex supports at depth $\ge R_0+12$ in $\Omega_2$; S3: scalar line ends at depth $\ge R_0$). At the uses these depths are supplied by the nesting of Ba{\l}aban's regions: the vertex localisations lie within $1$ of $\Omega_1$ or of $\supp\tilde A$ (\cite[p.~412]{Ba3}: the $A'$-legs carry $g_k$ with $\supp g_k\subset\Omega_1$; \cite[p.~414]{Ba3}: the block vertices \cite[(1.13)--(1.15)]{Ba3} are paired at a common $y$), $\dist(\supp\tilde A,\partial\Omega)>2r(t)\ge2R\ge R_0+2$ under $R\ge3R_0+15$, and at the domain changes of \cite[pp.~579--580]{Ba2} the vertex supports have depth $\ge5r(t)-9$ in the smaller domain. This supply is part of the reading of the step in \cref{app:a06_BALtorus}; here the depths are hypotheses.

\paragraph{Inputs.} Three propagator estimates are used, all of them statements about $G_k(X,A)$ and the fluctuation covariance at a single scale, proved in \cref{app:a06_BALtorus} (group (A) there, \cref{thm:a06A:groupA}) from \cite{Ba4}:
\begin{enumerate}[label=\textup{(\roman*)},nosep,leftmargin=*]
\item \textbf{(P)} \cite[Prop.~2.1]{Ba1} $=$ the theorem \cite[(1.9)--(1.12)]{Ba4}: for $\Omega=B^k(\Omega^{(k)})$ a union of large blocks and $A$ regular in the sense of \cite[(2.23)]{Ba1}, $|(D^\eta_{A,\mu}G_k(\Omega,A)f)(x)|$ and $|(G_k(\Omega,A)f)(x)|$ are $\le c_0e^{-\delta_0\dist(x,\supp f)}\norm f_\infty$ for $\dist(x,\Omega^c)\ge R_0$ \cite[(2.25)]{Ba1}, and the H\"older quotient $|x-x'|^{-\alpha}|U(A(\Gamma_{x,x'}))(DGf)(x')-(DGf)(x)|$ obeys the same bound with $\dist(\{x,x'\},\supp f)$ in the exponent, for $\dist(\{x,x'\},\Omega^c)\ge R_0$ \cite[(2.24)]{Ba1}. The constants $\delta_0,c_0,R_0$ are independent of $A$, $k$, $\Omega$ and of the torus, and the depth restriction is omitted for $\Omega=T_\eta$ \cite[pp.~610--611]{Ba1}.
\item \textbf{(C)} \cite[Prop.~2.3, (2.33)--(2.34)]{Ba1} at $L=5$: $|C^{(k)}_\Lambda(X,A;y,y')|\le c_2e^{-\theta_0|y-y'|_\infty}$ with no restriction on the points, and $(C^{(k)}_\Lambda)^{-1}\ge\gamma_X>0$; the constants depend on $L$ (through $\gamma_L$) and not on $k$, $X$, $\eps$ or the torus.
\item \textbf{(H0)} \cite[(1.8)]{Ba4} on every union $X$ of unit blocks and on $T_\eta$: $H_X(A)-m^2t^2\ge\gamma_0\defeq\frac14\min\{2,a_*\}$, for $A$ regular, which \eqref{sup1:eq:E} supplies.
\end{enumerate}

\paragraph{The norm \cite[(1.32)]{Ba3} for kernels.} For a kernel $Y(x,x')$ ``extended in a natural way'' to two variables \cite[p.~420]{Ba3} we use the strongest natural reading: $\norm Y_{1,\alpha}$ is the maximum over $a,b\in\{0,1\}$ of the suprema of $\norm{Y_{ab}}$, $Y_{ab}\defeq D^aYD^{b*}$, and of the $\alpha$-H\"older quotients in $x$ (for $a=1$), in $x'$ (for $b=1$) and of the mixed second difference (for $ab=1$), with partners $0<|x-\tilde x|\le1$ and $\tilde B$-transport along shortest contours. Quotients over larger distances are at most twice a supremum, and the summed form of \cite[(1.32)]{Ba3} is at most $16$ times this maximum.
